
\def\PUDcodigo{ECA214}

\def\PUDcodacad{04507.23}

\def\PUDdisciplina{Eletrônica II}

\def\PUDsemestre{S5}

\def\PUDhoras{80}

%NOVOS ---------------------------------------------------
\def\PUDhorasTeoria{60}
\def\PUDhorasPratica{20}

\def\PUDinterECA{ 
\hyperref[PUD:ECA211]{Eletrônica I (04507.15)}

\hyperref[PUD:ECA210]{Circuitos Elétricos I (04507.20)}
}

\def\PUDatuaisECA{---}

\def\PUDrecursos{Projetor multimídia; Quadro branco e pincel; Aulas em laboratório de eletrônica.}

%----------------------------------------------------------

\def\PUDcreditos{4}

\def\PUDprerequisito{ \hyperref[PUD:ECA210]{ECA210} }

\def\PUDpreacad{ \hyperref[PUD:ECA210]{Circuitos Elétricos I (04507.20)} }

\def\PUDobjetivos{Compreender o funcionamento dos principais dispositivos semicondutores empregados nos equipamentos eletrônicos, resolver os principais problemas relacionados aos circuitos eletrônicos, compreender o funcionamento e montar amplificadores básicos a BJT, com amplificadores operacionais (Amp-ops) e empregando o Transistor a Efeito de Campo (FET).}

\def\PUDementa{Introdução;  Circuitos com diodos;  Diodos para aplicações Especiais;  Transistores de junção bipolar (BJT);  Polarização do transistor Bipolar;  Modelos CA dos transistores bipolares;  Circuitos com transistores.  Amplificadores de potência empregando BJT;  Amplificadores operacionais;  Transistores de Efeito de campo (FET).}

\def\PUDconteudo{ & DIODO RETIFICADOR: Junção P-N e o Diodo ideal.  simbologia.  polarização.  capacitâncias.  Tensão de ruptura e corrente direta máxima.  curva característica.  circuitos com diodos. 
\\ 
 & CIRCUITOS COM DIODOS: Aproximações para o diodo.  Reta de carga.  Técnicas de análise de circuitos com diodos. 
\\ 
 & DIODOS PARA APLICAÇÕES ESPECIAIS: Modelo do Diodo Zener.  Máxima e mínima corrente de Zener.  Regulador de tensão a Zener.  Diodos Schottky.  Diodos Emissores de Luz – LED, Modelos.  Principais parâmetros dos LEDs. Principais parâmetros dos LEDs. Tecnologia atual de LEDs: Brancos, Azuis e outras cores.
\\ 
 & TRANSISTORES DE JUNÇÃO BIPOLAR: Transistor polarizado.  correntes no transistor.  Curvas de base e de coletor.  Folhas de dados dos transistores.  Reta de carga para o BJT.  Identificando a saturação.  Transistor como interruptor ou como fonte de corrente; 
\\ 
 & POLARIZAÇÃO DO TRANSISTOR BIPOLAR: Polarização com corrente de base constante.  Polarização com corrente de emissor constante.  Polarização combinada.  Polarização com tensão de coletor constante. 
\\ 
 & MODELOS CA DOS TRANSISTORES BIPOLARES: Modelo de Ebers-Moll.  Modelo PI do transistor.  Amplificadores com polarização de base.  Polarização de emissor e de coletor.  Aplicação dos transistores em Amplificadores de tensão. 
\\ 
 & AMPLIFICADORES DE POTÊNCIA EMPREGANDO BJT: Buffers de tensão.  Fontes lineares controladas e reguladas de tensão e corrente. 
\\ 
 & AMPLIFICADORES OPERACIONAIS: Amp-ops básicos.  Circuitos amp-ops práticos.  Especificações do amp-op.  Folha de dados do Amp-op.  Aplicações do amp-op. 
\\ 
 & TRANSISTORES DE EFEITO DE CAMPO: Operação do FET e IGFET.  Circuitos de polarização do FET.  Circuitos de polarização do IGFET.  Comportamento com a temperatura.  Amplificadores a FET: Fonte comum, dreno comum, gate (porta) comum.  Considerações para altas freqüências. }

\def\PUDmetodo{Aulas expositórias que podem ser teóricas e/ou práticas, onde as práticas no laboratório serão marcadas no decorrer da disciplina. No decorrer da disciplina é realizada a análise de circuitos eletrônicos e diagramas esquemáticos empregados industrialmente: circuitos para comunicação serial, filtros, temporizadores, aplicações em audio e em sensores industriais.}

\def\PUDavalia{A avaliação será feita com aplicação de provas teóricas e/ou práticas, além da possibilidade de inclusão de trabalhos e seminários no decorrer da disciplina.}

\def\PUDbibbasica{ & \href{www.biblioteca.ifce.edu.br/index.asp?codigo_sophia=24406}{Boylestad}, Robert L.; Nashelsky, Louis.  Dispositivos Eletrônicos e Teoria de Circuitos.  8ª ed.  São Paulo, SP: Pearson Prentice Hall, 2004.  672p.  ISBN: 9788587918222.
\\ 
 & \href{www.biblioteca.ifce.edu.br/index.asp?codigo_sophia=23875}{Malvino}, Albert Paul.  Eletrônica – Volume 1.  4ª ed.  São Paulo, SP: Pearson Makron Books, 1997.  670p.  ISBN: 9788534603782.
\\ 
 & \href{www.biblioteca.ifce.edu.br/index.asp?codigo_sophia=23930}{Malvino}, Albert Paul.  Eletrônica – Volume 2.  4ª ed.  São Paulo, SP: Pearson Makron Books, 1997.  558p.  ISBN: 853460455-X. }

\def\PUDbibcomplementar{ & \href{www.biblioteca.ifce.edu.br/index.asp?codigo_sophia=24345}{Smith}, Kenneth C.; Sedra, Adel S.  Microeletrônica.  5ª ed.  São Paulo, SP: Pearson Prentice Hall, 2007.  848p.  ISBN: 9788576050223.
\\ 
 & \href{www.biblioteca.ifce.edu.br/index.asp?codigo_sophia=24478}{Cruz}, Eduardo Cesar Alves.  Eletrônica aplicada.  2ª ed.  São Paulo, SP: Érica, 2012.  295p.  ISBN: 9788536501505.
\\ 
 & \href{www.biblioteca.ifce.edu.br/index.asp?codigo_sophia=24459}{Cipelli}, Antonio Marco V.  Teoria e desenvolvimento de projetos de circuitos eletrônicos.  23º ed.  São Paulo, SP: Érica, 2012.  445p.  ISBN: 9788571947597.
\\
 %& \href{www.biblioteca.ifce.edu.br/index.asp?codigo_sophia=56373}{Ahmed}, Ashfaq.  Eletrônica de Potência.  E-book.  Pearson.  484p.  ISBN: 9788587918031.
%\\
 & \href{www.biblioteca.ifce.edu.br/index.asp?codigo_sophia=23685}{Urbanetz Júnior}, Jair.  Eletrônica aplicada.  Curitiba, PR: Base Editorial, 2010.  144p.  ISBN: 9788579055751. }

\def\PUDautor{Luiz Daniel Santos Bezerra}

\def\PUDdata{2013-04-23}

\def\PUDrevisao{2}

\def\PUDrevisor{Luiz Daniel Santos Bezerra}

\def\PUDdatarevisao{2019-05-15}

\PUDcorpo
